% Figure: schematic Moody diagram showing the friction-factor regimes
% in pipe flow, with the worked Re = 1e6, e/D = 0.00015 point.
\begin{figure}[htbp]
\centering
\begin{tikzpicture}
\begin{loglogaxis}[
  width=12cm, height=7.5cm,
  xlabel={Reynolds number $\mathrm{Re} = \rho v D / \mu$},
  ylabel={Darcy friction factor $f$},
  xmin=5e2, xmax=1e8,
  ymin=0.008, ymax=0.1,
  grid=both,
  major grid style={engblack!20},
  minor grid style={engblack!10},
  legend pos=north east,
  legend cell align=left,
  font=\small,
]

% Laminar regime: f = 64/Re for Re < ~2300
\addplot[engblue, very thick, domain=5e2:2300, samples=20]
  {64/x};
\addlegendentry{Laminar: $f = 64/\mathrm{Re}$};

% Smooth-pipe turbulent (Blasius for moderate Re, Prandtl-von Karman
% asymptote for high Re). Use Blasius f = 0.316 Re^{-1/4} for
% 4000 <= Re <= 1e5, then a fitted decay for higher Re.
\addplot[engred, very thick, domain=4000:1e5, samples=80]
  {0.316/(x^0.25)};
\addlegendentry{Smooth: Blasius $f = 0.316\,\mathrm{Re}^{-1/4}$};

% Smooth-pipe extension (schematic): use a slower decay above 1e5.
\addplot[engred, very thick, dashed, domain=1e5:1e8, samples=80]
  {0.316/(x^0.25) * (x/1e5)^(-0.05)};

% Fully turbulent (rough pipe) horizontal asymptotes for several e/D.
% e/D = 0.05: f -> 0.072
\addplot[engblack, thick, dotted, domain=1e5:1e8] {0.072};
\node[engblack, font=\scriptsize] at (axis cs:1e7, 0.078)
  {$e/D = 0.05$};

% e/D = 0.005: f -> 0.031
\addplot[engblack, thick, dotted, domain=2e5:1e8] {0.031};
\node[engblack, font=\scriptsize] at (axis cs:1e7, 0.034)
  {$e/D = 0.005$};

% e/D = 0.00015 (commercial steel): f -> 0.013 at very high Re
\addplot[engblack, thick, dotted, domain=4e5:1e8] {0.013};
\node[engblack, font=\scriptsize] at (axis cs:1e7, 0.0115)
  {$e/D = 0.00015$};

% Mark the worked point: Re = 1e6, e/D = 0.00015, f approx 0.014
\addplot[engred, mark=*, mark size=2pt, only marks]
  coordinates {(1e6, 0.014)};
\node[engred, font=\scriptsize, anchor=south] at (axis cs:1e6, 0.015)
  {$f \approx 0.014$};

% Transition region label
\node[engblack, font=\scriptsize, align=center]
  at (axis cs:3000, 0.06)
  {transition\\$2300 < \mathrm{Re} < 4000$};

% Regime labels
\node[engblue, font=\scriptsize] at (axis cs:1500, 0.025)
  {laminar};
\node[engred, font=\scriptsize] at (axis cs:2e4, 0.022)
  {smooth turbulent};
\node[engblack, font=\scriptsize] at (axis cs:5e6, 0.06)
  {fully rough};

\end{loglogaxis}
\end{tikzpicture}
\caption[Moody diagram, schematic]{Schematic Moody diagram. Below
$\mathrm{Re} \approx 2300$ the friction factor obeys
$f = 64/\mathrm{Re}$ (Hagen-Poiseuille). Between $\mathrm{Re}
\approx 2300$ and $\mathrm{Re} \approx 4000$ the flow is in
transition; the friction factor is unstable and a designer treats
this band as unusable. For $\mathrm{Re} > 4000$, the friction factor
depends on both $\mathrm{Re}$ and the relative roughness $e/D$ per
the Colebrook equation \cite{paper:colebrook-1939}; at high
$\mathrm{Re}$ the dependence on $\mathrm{Re}$ disappears and $f$
asymptotes to a value set by $e/D$ alone (fully rough regime). The
marked point shows a $6\,\text{in}$ schedule-40 commercial steel
pipe ($e \approx 46\,\si{\micro\meter}$, $D \approx 154\,\si{\milli
\meter}$, $e/D \approx 3 \times 10^{-4}$) at $\mathrm{Re} = 10^{6}$,
giving $f \approx 0.014$ on the Swamee-Jain explicit form of
Colebrook \cite{paper:swamee-jain-1976}.}
\label{fig:vol01:ch02:moody}
\end{figure}
